\section{Dinâmica de Atitude}
A atitude de uma espaçonave livre de restrições externas pode ser descrita pelas equações de Euler para corpo rígido.
Considerando o sistema em torno do centro de massa, as equações de movimento rotacional são dadas por:
\begin{equation}
\mathbf{I}\,\dot{\vec{\omega}} + \vec{\omega} \times \big(\mathbf{I}\,\vec{\omega}\big) = \vec{M},
\label{eq:euler_rigidbody}
\end{equation}

onde $\mathbf{I}$ é a matriz de inércia em relação ao centro de massa, $\vec{\omega}$ é o vetor de velocidade angular expresso no sistema de corpo, e $\vec{M}$ corresponde ao vetor de momentos aplicados.

\subsection{Matriz de inércia e momentos aplicados}
A matriz de inércia da base não é constante quando existe um manipulador acoplado, uma vez que a movimentação das juntas altera a distribuição de massa do sistema.
Sendo assim, a matriz equivalente pode ser escrita como
\begin{equation}
\mathbf{I}(t) = \mathbf{I}_{\text{base}} + \mathbf{I}_{\text{manip}}(t),
\label{eq:inercia_total}
\end{equation}

sendo $\mathbf{I}_{\text{base}}$ a matriz fixa da base principal e $\mathbf{I}_{\text{manip}}(t)$ a contribuição variável proveniente do manipulador.  
Os momentos aplicados incluem tanto os torques de controle, quanto os torques de reação oriundos do movimento das juntas.

\subsection{Efeito do manipulador robótico acoplado à base}
O movimento do manipulador robótico provoca variação inercial e gera acoplamento dinâmico entre a base e os elos móveis.
A dinâmica acoplada pode ser representada por:
\begin{equation}
\mathbf{I}(t)\,\dot{\vec{\omega}} + \vec{\omega} \times \big(\mathbf{I}(t)\,\vec{\omega}\big) = \vec{M}_{\text{controle}} + \vec{M}_{\text{manip}},
\label{eq:dinamica_acoplada}
\end{equation}

onde $\vec{M}_{\text{controle}}$ são os torques de controle aplicados à espaçonave e $\vec{M}_{\text{manip}}$ representam os torques de reação devido ao movimento das juntas do manipulador.

Assim, a modelagem da dinâmica de atitude requer a consideração combinada da matriz de inércia variável, dos torques de controle e dos efeitos do manipulador robótico.
